\section{Введение}

\subsection{Актуальность темы}
Выпуклые множества являются важным объектом исследования в математике,
геометрии и методах оптимизации. Многие практические задачи формулируются
с использованием выпуклых областей допустимых решений и выпуклых
ограничений. Поэтому изучение свойств выпуклости позволяет лучше понимать
структуру оптимизационных задач и выбирать подходящие методы их решения.


\subsection{Цель работы}
Целью лабораторной работы является экспериментальное исследование выпуклых множеств и их свойств с использованием языка Python и библиотек NumPy, SciPy и Matplotlib.

В работе решаются следующие задачи:
\begin{itemize}
	\item построение выпуклой оболочки случайного набора точек на плоскости;
	\item визуальная проверка выпуклости пересечения множеств;
	\item программная проверка выпуклости заданного многоугольника.
	\item визуализация полученных результатов;
    \item сравнение экспериментальных результатов с теоретическими
    свойствами выпуклых множеств.
\end{itemize}

Помимо реализации алгоритмов, необходимо сопоставить вычислительные
результаты с теоретическими свойствами выпуклых множеств и ответить на
контрольные вопросы из технического задания:


